% !TeX program = xelatex
% UTF-8 transcription of the three handwritten source pages.
% Values, arrows and terminology are transcribed, not medically revised.
% pic01.jpg: 20261013_064343400.jpg (source page 1)
% pic02.jpg: 20261013_064350995.jpg (source page 2)
% pic03.jpg: 20261013_064357476.jpg (source page 3)
% Compile standalone with XeLaTeX, or input into a XePersian document.
\ifdefined\settextfont
  \def\ParathyroidStandalone{0}
\else
  \def\ParathyroidStandalone{1}
  \documentclass[12pt,a4paper]{article}
  \usepackage[margin=2cm]{geometry}
  \usepackage{amsmath,amssymb}
  \usepackage{xepersian}
  \IfFontExistsTF{Amiri}{\settextfont{Amiri}}{\settextfont{FreeSerif}}
  \setlatintextfont{FreeSerif}
\fi
\providecommand{\en}[1]{\lr{#1}}
\newcommand{\ptarrow}{\en{$\longleftarrow$}}
\newcommand{\ptmark}[1]{\textbf{#1}}
\newcommand{\ptuncertain}[1]{#1\,\textsuperscript{[خوانش نامطمئن]}}
\newcommand{\ptpage}[1]{\clearpage\section*{اختلالات پاراتیروئید ـ صفحهٔ #1}}
\ifnum\ParathyroidStandalone=1 \begin{document}\fi

\ptpage{۱}
{\footnotesize عنوان دست‌نویس: ۱۷۴ ـ اختلالات پاراتیروئید، \ptuncertain{اندو}.\par}
\begin{itemize}
\item کاهش کلسیم یونیزهٔ سرم \ptarrow\ ترشح \en{PTH}.
\item \en{PTH} در کلیه \ptarrow\ بازجذب کلسیم، ترشح فسفات.
\item کمبود \en{PTH} در آزمایش \ptarrow\ \en{$\mathrm{Ca}\downarrow$}، فسفات \en{$\uparrow$}.
\end{itemize}

\subsection*{هایپوکلسمی}
\begin{itemize}
\item افزایش تحریک‌پذیری نوروموسکولار:
  \begin{itemize}
  \item کرامپ عضلانی، اسپاسم کارپوپدال، ضعف، پارستزی.
  \item لارنگواسپاسم، تشنج.
  \item \en{Chvostek sign} \ptarrow\ نشانهٔ شوستک؛ اسپاسم صورت.
  \item \en{Trousseau sign} \ptarrow\ اسپاسم کارپال پس از \en{3--5} دقیقه؛ \en{15 mmHg} بالاتر از فشار سیستولیک.
  \end{itemize}
\item کلسیم در سرم:
  \begin{itemize}
  \item یونیزه.
  \item متصل ـ آلبومین.
  \end{itemize}
\item هرچه \en{pH} بالاتر برود، اتصال کلسیم ـ آلبومین بیشتر می‌شود.
\item \ptmark{در آلکالوز \ptarrow\ کلسیم یونیزه \en{$\downarrow$} \ptarrow\ علائم هایپوکلسمی}.
\item هایپوپروتئینی \ptarrow\ کلسیم یونیزه \en{N/L} ولی توتال \en{$\downarrow$} است.
\end{itemize}

\subsection*{پرایمری هایپوپاراتیروئیدی}
\begin{itemize}
\item هایپوکلسمی رخ می‌دهد ولی راشیتیسم نه.
\item اتیولوژی پرایمری هایپوپارا:
  \begin{itemize}
  \item مالفورمیشن‌های مادرزادی \ptarrow\ \ptmark{سندرم دی‌جورج}.
  \item آسیب به دنبال جراحی.
  \item خودایمنی.
  \end{itemize}
\end{itemize}

\ptpage{۲}
\subsection*{سودوهایپوپاراتیروئیدی}
\noindent \en{$\mathrm{Ca}\downarrow$}، فسفات \en{$\uparrow$}.
\begin{itemize}
\item تیپ \en{Ia}: هورمون \en{PTH} زیاد است ولی نمی‌تواند گیرنده را تحریک کند.
  \begin{latin}
  \noindent $G_{s\alpha}$ protein $\times$, Adenylate cyclase $\checkmark$.
  \end{latin}
\item تیپ \en{Ib}: فنوتیپ \en{N/L}.
  \begin{latin}
  \noindent $G_{s\alpha}$ protein $\checkmark$, Adenylate cyclase $\times$.
  \end{latin}
\item تیپ \en{Ic}: فنوتیپ غیرطبیعی \ptarrow\ هر دو سالم ولی یک نقص دیگر جلوی اثر \en{PTH} را می‌گیرد.
\item تیپ \en{II}: فنوتیپ \en{N/L} \ptarrow\ هر دو سالم، نقص پس از \ptuncertain{سنتز}.
  \begin{itemize}
  \item \ptuncertain{حاشیهٔ زیر شاخهٔ تیپ II: خیر}.
  \end{itemize}
\item سودوهایپوپارا \ptarrow\ وراثت \ptmark{\en{AD}}، تظاهر از تولد یا بعدتر.
\item با بیماری زیر همراهی دارد: \ptmark{\en{Albright hereditary osteodystrophy}}.
\end{itemize}
\begin{latin}
\begin{itemize}
\item Short stature, stocky body habitus, round facies.
\item \textbf{Short 4th and 5th metacarpals}.
\item Calcification of basal ganglia, SubQ calcification, developmental delay.
\end{itemize}
\end{latin}
% A short crossed-out line between types Ia and Ib is illegible.

\subsection*{هایپوکلسمی گذرای نوزادی (\en{TNH})}
\begin{itemize}
\item \ptmark{ظرف ۳ روز اول} پس از تولد به‌طور طبیعی \en{Ca} سرم کم می‌شود ولی در این بیماران به سرعت برنمی‌گردد.
\item نوزاد مادری که هایپوپاراتیروئید است \ptarrow\ هایپوکلسمیک.
\end{itemize}

\subsection*{منیزیم و ترشح \en{PTH}}
\begin{itemize}
\item برای ترشح \en{PTH} نیاز به منیزیم هست.
\item \ptmark{هایپومنیزیمی \ptarrow\ هایپوپارای ثانویه} \ptarrow\ درمان: جایگزینی منیزیم.
\item بهتر است \en{PTH} سرم زمانی اندازه‌گیری شود که \en{Ca} سرم پایین باشد.
\end{itemize}

\subsection*{درمان تتانی یا تشنج ناشی از هایپوکلسمی}
\begin{latin}
\noindent \textbf{IV calcium gluconate} (1--2 ml/kg of 10\% solution) over 10 minutes.
\end{latin}
\begin{itemize}
\item بررسی قلب با \en{EKG} از نظر \ptmark{برادی‌کاردی}.
\end{itemize}

\ptpage{۳}
\subsection*{درمان طولانی‌مدت}
\begin{itemize}
\item کلسیم + ویتامین \en{D (1,25)}.
\item سطح سرمی \en{Ca} \ptarrow\ \en{Lower half of N/L range}.
\end{itemize}

\subsection*{راشیتیسم}
\begin{itemize}
\item نقص در مینرالیزه‌شدن استخوان در کودکان.
\item در بزرگسالان به آن استئومالاسی می‌گویند.
\item استخوان نرم است.
\item \ptmark{متافیز استخوان \en{Wide} می‌شود}.
\item رشد قدی ناکافی، \ptmark{پای پرانتزی} (دردناک ممکن است باشد).
\item ضخیم‌شدن مچ دست و زانو.
\item در \en{X-ray}: متافیز دنده (کستوکندرال) باز شده است.
\end{itemize}

\subsection*{کمبود ویتامین \en{D}}
\begin{itemize}
\item کمبود ویتامین \en{D} تغذیه‌ای.
\item سوءجذب چربی می‌تواند عامل آن باشد:
  \begin{itemize}
  \item بیماری‌های هپاتوبیلیاری، \en{CF}.
  \end{itemize}
\item عامل دیگر: نقص در متابولیسم \en{vit D} در کلیه:
  \begin{itemize}
  \item نارسایی کلیه، کمبود \en{AR} آنزیم.
  \end{itemize}
\item نقص در متابولیسم \en{vit D} در کبد.
\end{itemize}

\subsection*{\en{Familial hypophosphatemic rickets}}
\begin{itemize}
\item \ptmark{کلیه نمی‌تواند فسفات را بازجذب کند}.
\item فسفات سرم \en{$\downarrow$}، فسفات ادرار \en{$\uparrow$}.
\item وراثت وابسته به \en{X}؛ چند ژن هم دارد.
\end{itemize}

\subsection*{بررسی اتیولوژی راشیتیسم}
\begin{itemize}
\item ارزیابی مواد معدنی، \en{vit D}.
\item کمبود تغذیه‌ای \ptarrow\ \ptmark{\en{25-hyd vit D $<8\,\mathrm{ng/ml}$}}.
\end{itemize}

\subsection*{درمان راشیتیسم}
\begin{itemize}
\item کلسیم + \en{vit D}.
\item در هایپوفسفاتمیک \ptarrow\ فسفات + \en{vit D}.
\end{itemize}

% Printed stationery: Subject / Year / Month / Date; chapter number 174
% on each page and handwritten source page numbers 1, 2, 3.
% Yellow highlighting and circled headings are represented by bold text
% and section headings. Blank ruled space has no content.
\ifnum\ParathyroidStandalone=1
  \let\ParathyroidEnd\enddocument
\else
  \let\ParathyroidEnd\relax
\fi
\ParathyroidEnd
